% SEGMENT OLP-0121-S01
% Part: first-order-logic
% Chapter: axiomatic-deduction
% Section: provability-consistency

\documentclass[../../../include/open-logic-section]{subfiles}

\begin{document}

\iftag{FOL}
      {\olfileid{fol}{axd}{prv}}
      {\olfileid{pl}{axd}{prv}}

\olsection{\usetoken{S}{derivability} மற்றும் முரணின்மை}

!!{derivability} உறவின் பல பண்புகளை இப்போது நிறுவுவோம். அவை தனித்தும்
சுவையானவை; மேலும் அவை ஒவ்வொன்றும் நிறைவுத் தேற்றத்தின் நிறுவலில் ஒரு பங்கு
வகிக்கும்.

% SEGMENT OLP-0121-S02
\begin{prop}\ollabel{prop:provability-contr}
  $\Gamma \Proves !A$ என்றும், $\Gamma \cup \{!A\}$ முரணுடையது என்றும்
  இருந்தால், $\Gamma$ முரணுடையது.
\end{prop}

\begin{proof}
  $\Gamma \cup \{!A\}$ முரணுடையது எனில், $\Gamma \cup \{!A\}
  \Proves \lfalse$. \olref[ptn]{prop:reflexivity} படி, ஒவ்வொரு $!B \in
  \Gamma$ க்கும் $\Gamma \Proves !B$. கருதுகோளின்படி $\Gamma \Proves !A$
  உம் என்பதால், ஒவ்வொரு $!B \in \Gamma \cup \{!A\}$ க்கும் $\Gamma
  \Proves !B$. \olref[ptn]{prop:transitivity} படி $\Gamma \Proves
  \lfalse$; அதாவது, $\Gamma$ முரணுடையது.
\end{proof}

% SEGMENT OLP-0121-S03
\begin{prop}
\ollabel{prop:prov-incons}
$\Gamma \Proves !A$ என்பது, அப்போதும் அப்போது மட்டுமே, $\Gamma \cup
\{\lnot !A\}$ முரணுடையது.
\end{prop}

\begin{proof}
முதலில் $\Gamma \Proves !A$ என்க. அப்போது \olref[ptn]{prop:monotonicity}
படி $\Gamma \cup \{\lnot !A\} \Proves !A$. மேலும்
\olref[ptn]{prop:reflexivity} படி $\Gamma \cup \{\lnot !A\} \Proves
\lnot !A$. \olref[prp]{ax:lnot2} படி $\Proves \lnot !A \lif (!A \lif
\lfalse)$ உம். ஆகவே \olref[ded]{prop:mp} ஐ இருமுறை பயன்படுத்தினால்,
$\Gamma \cup \{\lnot !A\} \Proves \lfalse$.

இப்போது $\Gamma \cup \{\lnot !A\}$ முரணுடையது என்க; அதாவது, $\Gamma
\cup \{\lnot !A\} \Proves \lfalse$. கழித்தல் தேற்றத்தால் $\Gamma
\Proves \lnot !A \lif \lfalse$. \olref[prp]{ax:lfalse2} படி $\Gamma
\Proves (\lnot !A \lif \lfalse) \lif \lnot\lnot !A$; ஆகவே
\olref[ded]{prop:mp} படி $\Gamma \Proves \lnot\lnot !A$. $\Gamma \Proves
\lnot\lnot !A \lif !A$ (\olref[prp]{ax:dne}) என்பதால்,
\olref[ded]{prop:mp} ஐ மீண்டும் பயன்படுத்தி $\Gamma \Proves !A$ பெறுகிறோம்.
\end{proof}

% SEGMENT OLP-0121-S04
\begin{prob}
$\Gamma \Proves \lnot !A$ என்பது, அப்போதும் அப்போது மட்டுமே, $\Gamma
\cup \{!A\}$ முரணுடையது என்பதை நிறுவுக.
\end{prob}

% SEGMENT OLP-0121-S05
\begin{prop}\ollabel{prop:explicit-inc}
  $\Gamma \Proves !A$ என்றும் $\lnot !A \in \Gamma$ என்றும் இருந்தால்,
  $\Gamma$ முரணுடையது.
\end{prop}

\begin{proof}
  \olref[prp]{ax:lnot2} படி $\Gamma \Proves \lnot !A \lif (!A \lif
  \lfalse)$. \olref[ded]{prop:mp} ஐ இருமுறை பயன்படுத்தினால் $\Gamma
  \Proves \lfalse$.
\end{proof}

% SEGMENT OLP-0121-S06
\begin{prop}\ollabel{prop:provability-exhaustive}
  $\Gamma \cup \{!A\}$, $\Gamma \cup \{\lnot !A\}$ ஆகிய இரண்டும்
  முரணுடையவை எனில், $\Gamma$ முரணுடையது.
\end{prop}

\begin{proof}
பயிற்சி.
\end{proof}

\tagprob{FOL}
\begin{prob}
  \olref[fol][axd][prv]{prop:provability-exhaustive} ஐ நிறுவுக.
\end{prob}
\tagendprob

\tagprob{notFOL}
\begin{prob}
  \olref[pl][axd][prv]{prop:provability-exhaustive} ஐ நிறுவுக.
\end{prob}
\tagendprob

\end{document}

